\chapter{Gauss's law}\label{ch:gauss}

Gauss's law is Coulomb's law wearing symmetry: the flux through any
closed surface counts only the charge inside. It turns three-page
integrals into three-line algebra --- when the symmetry is real.

\section{Flux: counting field lines}\label{sec:ga-flux}

\begin{equation}
  \Phi_E=\int\mathbf E\cdot d\mathbf A,\qquad
  \oint\mathbf E\cdot d\mathbf A=\frac{Q_{\rm enc}}{\varepsilon_0},
  \label{eq:gauss}
\end{equation}
open surface vs closed (figure~\ref{fig:flux}). Flux through a closed
surface is charge inside divided by $\varepsilon_0$ --- the field's
bookkeeping made literal, and the first Maxwell equation in the book.

\begin{figure}[htbp]\centering
\begin{tikzpicture}[scale=0.9,font=\scriptsize]
\draw[very thick] (0,0) circle (0.5) node{$dA$};
\draw[->,thick,accent] (0,0) -- (0,0.9) node[above]{$\hat n$};
\draw[->,thick,amber] (-0.4,-0.4) -- (0.2,0.2);
\node at (0,-1.2) {open: $\Phi=\int\mathbf E\cdot d\mathbf A$};
\draw[very thick,accent] (4,0) ellipse (1.3 and 0.9);
\draw[->,thick,accent] (4,0.9) -- (4,1.3);
\draw[->,thick,accent] (4,-0.9) -- (4,-1.3);
\draw[->,thick,accent] (5.3,0) -- (5.8,0);
\draw[->,thick,accent] (2.7,0) -- (2.2,0);
\fill[amber] (4,0) circle (0.05);
\node at (4,-1.8) {closed: $\oint\mathbf E\cdot d\mathbf A=Q_{\rm enc}/\varepsilon_0$};
\end{tikzpicture}

\caption{Flux through an open patch vs a closed surface. Only the closed
integral counts charge.}
\label{fig:flux}
\end{figure}

\section{From Coulomb to Gauss: the intuition}\label{sec:ga-deriv}

A point charge's flux through a surrounding sphere is
$kq/r^{2}\cdot4\pi r^{2}=q/\varepsilon_0$ --- radius cancels, so any
shape works (deformations keep line count). Superposition extends it to
any distribution, and charge \emph{outside} contributes zero net flux:
its lines enter and leave. Our code check makes it concrete below.

\section{Spherical symmetry}\label{sec:ga-sphere}

Spherical symmetry means $\mathbf E=E(r)\hat r$ is constant on the
Gaussian sphere: $\oint E\,dA=E(r)\cdot4\pi r^{2}$, so
$E(r)=Q_{\rm enc}(r)/4\pi\varepsilon_0r^{2}$.

\begin{example}[Uniform-volume-charge sphere, both regions]\label{ex:sphere}
A sphere of radius $R$ carries $Q$ uniformly in volume. Find $E(r)$
inside and out.
\par\smallskip
\textbf{Strategy.} Enclosed charge is the part below radius $r$.
\begin{solution}
Outside ($r>R$): $Q_{\rm enc}=Q$, so $E=kQ/r^{2}$ --- the shell theorem
of chapter~8 again.
Inside ($r<R$): $Q_{\rm enc}=Q(r/R)^{3}$, so
$E=kQr/R^{3}$ --- linear, zero at centre (figure~\ref{fig:er}).
The matching point at $r=R$ checks both formulas: $kQ/R^{2}$.
\checkmark
\end{solution}
\end{example}

\begin{figure}[htbp]\centering
\begin{tikzpicture}
\begin{axis}[width=8.5cm,height=4.8cm,xlabel={$r$},ylabel={$E(r)$},
xmin=0,xmax=3,ymin=0,ymax=1.5,font=\small]
\addplot[very thick,accent,domain=0:1] {x} node[pos=0.5,below,font=\scriptsize]{$\propto r$ (inside)};
\addplot[very thick,accent,domain=1:3] {1/x^2} node[pos=0.7,above,font=\scriptsize]{$\propto1/r^{2}$};
\addplot[mark=*] coordinates {(1,1)};
\addplot[very thick,amber,domain=0:3] {0.7};
\node[font=\scriptsize] at (axis cs:1.5,0.85) {sheet: const};
\end{axis}
\end{tikzpicture}

\caption{$E(r)$ for a uniform sphere (rises to the surface, then
$1/r^{2}$) against the constant sheet field.}
\label{fig:er}
\end{figure}

\section{Cylindrical symmetry}\label{sec:ga-cyl}

For a line or cylinder, $\mathbf E=E(r)\hat r$ on the curved wall gives
$E\cdot2\pi rL=Q_{\rm enc}/\varepsilon_0$, so
$E=\lambda/2\pi\varepsilon_0r$: $1/r$, not $1/r^{2}$.

\begin{example}[Coaxial cable]\label{ex:coax}
Inner conductor carries $+\lambda$, the outer $-\lambda$. Find the
field between them and outside the cable.
\par\smallskip
\textbf{Strategy.} Enclosed charge on a coaxial Gaussian cylinder.
\begin{solution}
Between ($a<r<b$, figure~\ref{fig:coax}): only the inner wire enclosed,
$E=\lambda/2\pi\varepsilon_0r$ radially out.
Outside: net enclosed zero --- $E=0$. The cable contains its own field,
which is exactly why engineers buy coax. \checkmark
\end{solution}
\end{example}

\begin{figure}[htbp]\centering
\begin{tikzpicture}[scale=1.0,font=\scriptsize]
\draw[fill=accent!15,very thick] (0,0) circle (0.4);
\draw[very thick] (0,0) circle (1.0);
\draw[dashed,very thick,amber] (0,0) circle (0.7);
\node at (0,0) {$+\lambda$};
\node at (0,1.3) {outer cylinder};
\node at (0,-1.4) {Gaussian: only inner enclosed};
\draw[->,thick,amber] (0.7,0) -- (0.7,0.55) node[right,pos=0.6]{$E=\lambda/2\pi\varepsilon_0r$};
\node[align=left] at (0,-1.4) {};
\end{tikzpicture}

\caption{The coaxial geometry: between the conductors the field is
$1/r$; outside it's zero by cancellation.}
\label{fig:coax}
\end{figure}

\section{Planar symmetry}\label{sec:ga-plane}

For an infinite sheet, flux escapes through both pillbox faces:
$2EA=\sigma A/\varepsilon_0$, so $E=\sigma/2\varepsilon_0$ ---
\emph{constant}, no distance at all (figure~\ref{fig:pill}). Two sheets
of $\pm\sigma$ superpose to $\sigma/\varepsilon_0$ between and zero
outside --- the parallel-plate capacitor's field, free of charge.

\begin{example}[Sheet vs conductor plate]\label{ex:sheet}
A large conducting plate carries total $\sigma$ on \emph{both} faces
combined. Find $E$ just outside it, and compare with an insulating
sheet of the same $\sigma$.
\par\smallskip
\textbf{Strategy.} Conductor: $E=0$ inside material, so all flux exits
one way at each face. Sheet: both ways.
\begin{solution}
Conductor: pillbox with one face inside the metal gets $E=0$ there, so
$EA=\sigma A/\varepsilon_0$ gives $E=\sigma/\varepsilon_0$ (here
$\sigma$ is per-face). Insulating sheet: $E=\sigma/2\varepsilon_0$
each side. For the same total charge per area the conductor's surface
field is twice the sheet's --- the missing-face bookkeeping again.
\checkmark
\end{solution}
\end{example}

\begin{figure}[htbp]\centering
\begin{tikzpicture}[scale=1.0,font=\scriptsize]
\draw[very thick,accent] (-1.6,0) -- (1.6,0) node[right]{$\sigma$};
\draw[dashed,very thick] (-0.5,-0.5) rectangle (0.5,0.5);
\draw[->,thick,amber] (0,0.15) -- (0,0.75) node[above]{$E$};
\draw[->,thick,amber] (0,-0.15) -- (0,-0.75) node[below]{$E$};
\node[align=left] at (0,-1.3) {pillbox straddles surface:\\$E=\sigma/2\varepsilon_0$ sheet,\\$E=\sigma/\varepsilon_0$ conductor};
\end{tikzpicture}

\caption{The pillbox argument: two faces share the sheet's flux.}
\label{fig:pill}
\end{figure}

\section{Conductors: where $E$ goes to zero}\label{sec:ga-cond}

In electrostatic equilibrium $E=0$ inside conductor material (any field
would move the free charge until it vanished), so the conductor is one
equipotential body, all excess charge lives on the surface, and
$E=\sigma/\varepsilon_0$ just outside --- twice the sheet's field
because the interior side contributes nothing
(figure~\ref{fig:cavity}).

\begin{figure}[htbp]\centering
\begin{tikzpicture}[scale=0.9,font=\scriptsize]
\draw[very thick,rounded corners=0.3] (-2,-1.3) rectangle (0.4,1.3);
\draw[fill=white,draw=accent] (-0.8,0) circle (0.55);
\node at (-0.8,0) {$+q$};
\node[align=left] at (-0.8,-1.9) {charge in cavity:\\$-q$ on inner skin, $+q$ on outer};
\node[align=left] at (-0.8,-1.9) {};
\draw[very thick,rounded corners=0.3] (2.6,-1.3) rectangle (5,1.3);
\draw[fill=white,draw=accent] (3.8,0) circle (0.55);
\node at (3.8,0) {empty};
\node[align=left] at (3.8,-1.9) {empty cavity:\\$E=0$ inside, shielded};
\end{tikzpicture}

\caption{Conductor cavities: a charge inside induces $-q$ on the inner
skin and $+q$ on the outer; an empty cavity stays field-free.}
\label{fig:cavity}
\end{figure}

\begin{example}[Charge in a conductor cavity]\label{ex:cavity}
A point charge $+q$ sits in a cavity inside a neutral conductor. Find
the induced surface charges and the field outside.
\par\smallskip
\textbf{Strategy.} A Gaussian surface in the conductor material has
$E=0$ everywhere on it.
\begin{solution}
$\Phi=0$ forces $Q_{\rm enc}=0$: the inner wall carries exactly $-q$.
Conservation of the conductor's neutrality puts $+q$ on the outer
surface, so the outside world sees a field as if a point charge $+q$
sat at the conductor --- the cavity's contents announced, position
hidden. \checkmark
\end{solution}
\end{example}

\begin{example}[Enclosed charge by inspection]\label{ex:encl}
A spherical insulating shell carries $+Q$; a point charge $-q$
($q<Q$) sits at its centre. Find $E$ just inside and just outside the
shell.
\par\smallskip
\textbf{Strategy.} $Q_{\rm enc}$ is everything inside your surface ---
shell charge counts only beyond its radius.
\begin{solution}
Inside ($r<R$): only the point, $E=-kq/r^{2}$ (inward).
Outside: $Q_{\rm enc}=Q-q$, so $E=k(Q-q)/r^{2}$ outward.
At the shell radius the field jumps by $kQ/R^{2}=\sigma/\varepsilon_0$
--- the surface-charge contribution, exactly as \S\ref{sec:ga-cond}
predicts. \checkmark
\end{solution}
\end{example}

\begin{lstlisting}[language=Python,caption={\texttt{code/ch10\_flux\_check.py} --- numeric flux through a cube; convergence and the outside-charge zero.},label={lst:flux}]
En = sum(q*c/r**3 * nv for c, nv in zip(p, nvec))
flux += En*dA
\end{lstlisting}

Output of a real run:
\begin{lstlisting}[numbers=none]
Flux of point charge q=1 through enclosing cube (theory 4*pi=12.566):
  n= 10: Phi=12.5972  err=0.25%
  n= 30: Phi=12.5698  err=0.03%
  n= 90: Phi=12.5668  err=0.00%
  n=270: Phi=12.5664  err=0.00%

Charge OUTSIDE cube (expect ~0): Phi=-0.0000
\end{lstlisting}

A plain Riemann sum on a cube --- no symmetry exploited --- converges
to $4\pi$ (i.e.\ $q/\varepsilon_0$ in these units) and returns $0.0000$
for a charge outside: Gauss's law with the mesh resolution as the only
approximation.

\section{In engineering practice: shielding is Gauss's law}\label{sec:ga-eng}

A Faraday cage is \S\ref{sec:ga-cond} deployed: charge redistributes
until the interior field dies, which is why coax shields signals,
aircraft survive lightning strikes, and your phone misbehaves in a lift.
High-voltage labs take the corollary --- $E$ peaks at sharp curvature ---
and shape everything round to keep $\sigma$ (and hence the surface
field) down.

\section{Common mistakes}\label{sec:ga-mistakes}

\begin{tripwire}[Gauss applied without symmetry]
The law is always true but only \emph{useful} when $E$ is constant on
the surface: no symmetry, no shortcut --- go back to Coulomb
integration.
\end{tripwire}
\begin{tripwire}[Wrong enclosed charge]
$Q_{\rm enc}$ counts only charge \emph{inside} the Gaussian surface ---
a charge on the surface itself splits ambiguity, and outside charges
contribute nothing.
\end{tripwire}
\begin{tripwire}[Gaussian area vs charge area]
The $4\pi r^{2}$ is the \emph{Gaussian} sphere's area at radius $r$,
not the charged sphere's: use your surface's radius.
\end{tripwire}
\begin{tripwire}[$E=0$ inside conductor vs cavity]
Inside the \emph{material}: always zero. Inside an empty \emph{cavity}:
zero. Inside a cavity containing charge: emphatically not.
\end{tripwire}
\begin{tripwire}[Sheet vs conductor surface field]
Sheet: $\sigma/2\varepsilon_0$; conductor surface:
$\sigma/\varepsilon_0$. The factor of two is the missing second face ---
they are different setups, not a typo.
\end{tripwire}

\section{Summary}\label{sec:ga-summary}

\begin{table}[ht]\centering
\begin{tabularx}{\linewidth}{@{}lY@{}}
\toprule
\tablehead\thead{Idea} & \thead{One-line version}\\
\midrule
flux & $\Phi=\int\mathbf E\cdot d\mathbf A$; closed counts $Q_{\rm enc}$ \\
sphere & $E=kQ_{\rm enc}(r)/r^{2}$; inside uniform $\propto r$ \\
line & $E=\lambda/2\pi\varepsilon_0r$ \\
sheet & $E=\sigma/2\varepsilon_0$, constant \\
conductor & $E=0$ inside; $E=\sigma/\varepsilon_0$ at surface; shields cavities \\
\bottomrule
\end{tabularx}
\caption{Gauss's law cheat sheet.}
\label{tab:gasummary}
\end{table}

\begin{wins}[Definition of done]
You can now: pick the right Gaussian surface for sphere/line/sheet,
compute $E(r)$ in both regions, reason about enclosed charge, and argue
conductor/cavity induction from $E=0$.
\end{wins}

\chapternotes{%
The symmetry-first treatment follows MIT 8.02SC modules 5 and 8
\cite{ocw801} and AP Physics C E\&M \cite{apced}; Griffiths ch.~2 is the
deeper anchor (in print). The flux numbers are the listing's real
output.}

\begin{rulebox}[{Watch instead}]
\begin{enumerate}
\item \emph{Electric Flux, Gauss' Law, Examples} --- Walter Lewin, MIT
8.02x Lecture 3, 50:11, \url{https://youtu.be/Zu2gomaDqnM} ---
\S\ref{sec:ga-flux}--\S\ref{sec:ga-plane} at full length (accessed
2026-10-11).
\item \emph{Gauss law of electricity} --- Khan Academy India, 15:51,
\url{https://youtu.be/owMznEZPNt4} --- \S\ref{sec:ga-deriv} worked
quickly (accessed 2026-10-11).
\end{enumerate}
\end{rulebox}

\section*{Further reading}
\addcontentsline{toc}{section}{Further reading}

\begin{itemize}
\item MIT OCW \emph{8.02SC}, module 5 (Gauss's law) ---
recitations on all three symmetries. At \url{https://ocw.mit.edu},
course ``8.02SC'' (verified 2026-10-11).
\item OpenStax \emph{College Physics 2e}, ch.~18 ---
free text: \url{https://openstax.org} (verified 2026-10-11).
\item College Board, \emph{AP Physics C (E\&M) CED}, Unit 8 ---
verified 2026-10-11 at
\url{https://apstudents.collegeboard.org}.
\item E.~M. Purcell and D.~J. Morin, \emph{Electricity and Magnetism},
chs.~1--2 --- the deeper Gauss treatment.
\item R.~P. Feynman, \emph{Lectures on Physics}, vol.~II, chs.~4--5 ---
field intuition beyond the formulas; free at Caltech's site.
\item Khan Academy, \emph{Electric flux and Gauss's law} --- drills, at
\url{https://www.khanacademy.org} (verified 2026-10-11).
\end{itemize}

\section*{Additional exercises}
\addcontentsline{toc}{section}{Additional exercises}

\begin{exercise}[easy]
Compute $\Phi_E$ through a square face in a uniform field of
$300$\,N/C at $30^{\circ}$ to the normal. (AP-C)
\end{exercise}
\begin{exercise}[easy]
A cube contains charge $-2$\,\textmu C: find total flux and the sign
of the field's average direction. (AP-C)
\end{exercise}
\begin{exercise}[easy]
Why is $E=0$ inside a conductor in equilibrium --- and which step fails
if charge keeps moving? (AP-C)
\end{exercise}
\begin{exercise}[easy]
An infinite sheet carries $\sigma$: find $E$ both sides, then two
parallel sheets $\pm\sigma$. (F=ma)
\end{exercise}
\begin{exercise}[easy]
A sphere's surface charge is $\sigma$: find $E$ just outside and just
inside the surface. (AP-C)
\end{exercise}
\begin{exercise}[medium]
Re-derive the uniform sphere's $E(r)$ both regions; check the two
formulas agree at $r=R$. (AP-C FR-style)
\end{exercise}
\begin{exercise}[medium]
A line of charge $\lambda$: find $E(r)$ by Gauss and where the result
fails for a finite rod. (AP-C)
\end{exercise}
\begin{exercise}[medium]
A point charge sits at the centre of a spherical conducting shell:
find the charge on each surface and $E$ in all three regions. (F=ma)
\end{exercise}
\begin{exercise}[medium]
A coaxial cable has inner radius $a$, outer $b$: find where $E$ is
largest between conductors and why that limits the voltage.
(engineering)
\end{exercise}
\begin{exercise}[medium]
Two concentric shells carry $+Q$ and $-Q$: find $E(r)$ in all three
regions. (AP-C)
\end{exercise}
\begin{exercise}[hard]
A cavity inside a conductor holds charge $+q$ off-centre: argue the
outer surface charge is uniform and find the inner-wall total.
(USAPhO-style)
\end{exercise}
\begin{exercise}[hard]
Given $\rho(r)=\rho_0(1-r/R)$ inside a sphere, find $E(r)$ inside and
the charge distribution that would make $E$ constant inside.
(APhO-style)
\end{exercise}
\begin{exercise}[medium]
A flat disk carries $\sigma$: explain why the field on axis falls with
distance even though an infinite sheet's does not. (concept)
\end{exercise}
\begin{exercise}[medium]
A Gaussian cube replaces the sphere around a point charge: the law
still holds --- why is it useless for computing $E$? (concept)
\end{exercise}
\begin{exercise}[medium]
Find $E$ inside an empty cavity cut off-centre inside a uniformly
charged sphere. (F=ma; hint: superpose)
\end{exercise}
